\section{The Threshold}
\label{sec:ch01-the-threshold}

\index{adoption threshold|(}%
I went looking for a published threshold and did not find one. No engineer I
could find has written down a number of teams, or of services, above which
federation pays. The closest thing is a vendor guide with no author's name on
it, and I am not willing to build a chapter on a claim nobody signed. So what
follows is my judgment. I would rather say that plainly than dress it up as a
finding.

Four conditions. Federation pays when all four hold, and the count of your
teams is not one of them.

\begin{enumerate}[topsep=4pt, itemsep=3pt]
  \item More than one team needs to change the same API surface, and they
        cannot wait on each other to do it.
  \item Those teams already release independently. Separate pipelines,
        separate on-call, separate ability to roll back without a meeting.
  \item Every domain you intend to put in the graph has an owner who is
        allowed to say no to a schema change.
  \item Somebody will own the router and the composition pipeline as their
        product, not as a rota.
\end{enumerate}

The failure modes are more useful than the conditions. If the first is false,
you have one team, and you are about to solve a scheduling problem with a
distributed system. If the second is false, you get the entire cost of
federation and none of its benefit: the schema is split, and it still ships on
one train, so you have added a composer and a router to a release process that
was already the bottleneck. If the third is false, the composer will make the
ownership decision for you, at merge time, in the form of an error nobody feels
responsible for. If the fourth is false, the router is everyone's problem,
which means it is nobody's until it is down.

That second condition is the one I would test hardest, because it is the one
teams assume they pass. Deploying services separately is not the same as being
able to. If a schema change to your subgraph still needs a coordinated release
with two others, splitting the graph has changed the file layout and nothing
else.
\index{adoption threshold|)}%

\subsection{What Wins Below the Line}
\label{sec:ch01-what-wins-below-the-line}

\index{modular monolith}%
One GraphQL service with real module boundaries inside it. That is not a
consolation prize, and the argument for it does not depend on GraphQL at all.
Martin Fowler put it as a rule about starting: \enquote{you shouldn't start a
new project with microservices, even if you're sure your application will be
big enough to make it worthwhile}~\autocite{fowler2015}. His evidence was that
nearly every successful microservice system he had seen grew out of a monolith
that got too big, while the ones built distributed from the start tended to end
in trouble.

Shopify says the same thing from the other end, as a standing policy rather
than a starting rule. \enquote{We are very deliberate about when to split
functionality out into separate services, and we only do it for good reasons},
writes Philip M\uml{u}ller, and \enquote{That's because splitting a single
monolithic application into a distributed system of services increases the
overall complexity considerably}~\autocite{mueller2020}. That is a claim about
distribution generally, not about graphs, and it applies here because a
federated graph is a distributed system wearing a schema.

The GraphQL-specific version of the same choice runs at a scale most readers
will never need. Aditya Pratap Singh describes Zalando running one GraphQL
service across all of its domains: \enquote{instead of having multiple Graphs connected
via a library and gateway we have a single service at Zalando which connects
all the domains in a single schema Graph}, a schema that had grown \enquote{to
a dense graph now with more than 12 domains and serves more than 80\% of Web
and 50\% of the App use cases}~\autocite{singh2021}. More than twelve domains
in one service, covering most of that company's web use cases.

I want to be careful about what that last one proves, because it is easy to
over-read. Singh's post describes what Zalando built. It does not say they
evaluated federation and turned it down, and I am not going to imply that it
does. It establishes one thing only, and one thing is enough: a single graph
service is not a toy architecture you outgrow at three domains.
